\section{Direct estimates on finite pieces of the actual orbit}
\label{sec:direct-orbit-bounds}
A small residual for a compressed vector and a small error in a
finite-time compression are different requirements. We first record
their scales. We then obtain an estimate on the uncompressed dynamics
by a different argument: short-lag Gaussian estimates bound a finite
Gram matrix of exact orbit vectors. The resulting estimate holds on
every translate of that orbit segment. It involves no reconstruction
of a long-time vector on a spatial grid.

Throughout this and the next section the Hilbert space is
$\mathcal H_R=L^2(Y_R^*,\nu_R^*)$, with
$\langle f,g\rangle_R=\int\overline f g\dd\nu_R^*$.
Write
\begin{equation}\label{eq:direct-orbit-convention}
 \mathcal U_z=\mathcal U_{R,z},\qquad e=\one,\qquad
 c_{j,R}(z)=\langle e,\mathcal U_z^j e\rangle_R
       =\int e^{-iz\cdot U_{j,R}}\dd\nu_R^*,\quad j\ge0.
\end{equation}
Here $U_{j,R}=J_{j,R}-j\bar G_R$ is the centered record and
$\mathcal U_z f=e^{-iz\cdot g_R}f\circ F_R^*$ is the unitary
operator of \eqref{eq:actual-compressed-operator}. For negative $j$
we use its inverse powers; in particular
$c_{-j,R}(z)=\overline{c_{j,R}(z)}$. The symbol $z\in\R^4$ denotes
physical frequency throughout.

\subsection{The two grid scales}
\begin{proposition}[Separation of the displayed grid budgets]
\label{prop:grid-scale-separation}
Let $C_{\mathrm g}>0$ be the exponential variation constant in
\eqref{eq:exponential-compression-consistency}. Suppose $L_n\ge n$
and the grid contribution in that bound tends to zero, that is,
\begin{equation}\label{eq:grid-term-vanishing}
 n\sqrt{h_n(1+|z_n|)e^{C_{\mathrm g}L_n}}\longrightarrow0.
\end{equation}
If $|z_n|\ge2n^{-99/200}$ and
$\rho_n=(|z_n|^2+\xi_n^2)^{1/2}\le\rho_0<1$, then
\[
 \rho_n^2\{1+\log(C_0/(h_n\rho_n))\}\longrightarrow\infty.
\]
Thus these choices do not satisfy the local resolvent domain in
Corollary~\ref{cor:single-log-compressed}.
\end{proposition}
\begin{proof}
Squaring \eqref{eq:grid-term-vanishing} and taking logarithms gives
\[
 \log(1/h_n)-C_{\mathrm g}L_n-2\log n
                 -\log(1+|z_n|)\longrightarrow+\infty.
\]
Consequently $1+\log(C_0/(h_n\rho_n))\ge c n$ for all sufficiently
large $n$, with $c>0$. Since $\rho_n^2\ge4n^{-198/200}$, the product
is at least $4c n^{1/100}$. This proves the assertion.
\end{proof}

This statement concerns simultaneous use of the two displayed bounds,
not the size of the actual approximation error for every mesh. Neither
finite-dimensional theorem is contradicted. The direct argument below
makes no mesh choice and proves a different long-time conclusion:
average square cancellation, rather than decay at every individual
return count. The fixed-return raw inversion problem is unchanged.

\subsection{A uniform estimate on the short lags}
Put
\begin{equation}\label{eq:direct-exponent-definitions}
 \gamma=\frac1{28}-\frac1{200}=\frac{43}{1400},\qquad
 \alpha=\frac{200}{99},\qquad \kappa=\alpha-2=\frac2{99}.
\end{equation}

\begin{lemma}[Pointwise Gaussian comparison on the growing band]
\label{lem:pointwise-lag-Gaussian}
There is $C$ such that for $j\ge2$ and $|z|\le2j^{-99/200}$,
\begin{equation}\label{eq:pointwise-lag-Gaussian}
 \left|c_{j,R}(z)-e^{-jz^{\mathsf T}D_Rz/2}\right|
             \le Cj^{-\gamma}\sqrt{\log(2+j)}
\end{equation}
uniformly in $R$. This is a pointwise estimate, not an inference from
an integrated characteristic-function bound.
\end{lemma}
\begin{proof}
Set $v=-\sqrt j\,z$ and $\delta=j^{-1/14}/4$. Apply
Lemma~\ref{lem:explicit-frequency-error} at collision time
$\lfloor j/c_*\rfloor$, with initial density $s_R$, and use the
pointwise stopping estimate \eqref{eq:weighted-stopping-error}.
The rescaled frequency obeys $|v|\le2j^{1/200}$. Before integration,
the seven polynomial margins from smoothing, covariance replacement,
amplitude, cubic remainder, stopping and integer rounding are bounded
below by the entries in
\[
 \frac1{14},\quad\frac1{28}-\frac1{200},\quad
 \frac1{14}-\frac2{200},\quad\frac3{14}-\frac1{200},\quad
 \frac1{14}-\frac3{200},\quad\frac15-\frac1{200},\quad
 1-\frac2{200}.
\]
The second entry is the unique smallest one. It carries
$\sqrt{\log(2+j)}$; all other logarithmic factors are absorbed by
the strict larger margins. The complementary spectral power is
exponentially small. The analytic conditions hold because
$6/14+3/200<1/2$ and $2/14+1/200<1/2$. Enlarging the constant treats
the finite initial range of $j$, since the two transforms are bounded
by one. This proves \eqref{eq:pointwise-lag-Gaussian} directly from the
pointwise collision estimate.
\end{proof}

For $0<r=|z|\le r_0$, where $r_0$ is sufficiently small, choose
\begin{equation}\label{eq:direct-memory-scale}
 m=m(r)=\lfloor r^{-\alpha}\rfloor\ge2.
\end{equation}
Only lags shorter than $m$ will be estimated by Gaussian comparison.

\begin{proposition}[Summable short-lag budget]
\label{prop:short-lag-row-budget}
Uniformly in $R$ and $0<|z|=r\le r_0$,
\begin{equation}\label{eq:short-lag-row-budget}
 \mathcal B_m(R,z):=1+2\sum_{j=1}^{m-1}|c_{j,R}(z)|
  \le C\{r^{-2}+m^{1-\gamma}\sqrt{\log(2+m)}\},
 \qquad \frac{\mathcal B_m(R,z)}m\le Cr^\kappa.
\end{equation}
\end{proposition}
\begin{proof}
For every $2\le j<m$, the choice of $m$ gives
$rj^{99/200}\le1$, hence the physical frequency is in the band of
Lemma~\ref{lem:pointwise-lag-Gaussian}. Uniform ellipticity gives
$e^{-jz^{\mathsf T}D_Rz/2}\le e^{-d_0jr^2/2}$.
Summing the geometric series costs $Cr^{-2}$, and summing the error
costs $Cm^{1-\gamma}\sqrt{\log(2+m)}$. The lag one and diagonal
terms are absorbed in the first term. For small $r$ one has
$m\ge r^{-\alpha}/2$, so the normalized terms are at most
\[
 Cr^{\alpha-2},\qquad
 Cr^{\alpha\gamma}\sqrt{1+|\log r|}.
\]
Here $\alpha\gamma=43/693$ and $\kappa=14/693$; their difference is
$29/693>0$. The logarithm is therefore absorbed uniformly in
$Cr^\kappa$. No estimate on $c_j$ for $j\ge m$ is used.
\end{proof}

\subsection{Finite orbit Gram matrices}
\begin{lemma}[Analysis and synthesis of an exact orbit segment]
\label{lem:finite-orbit-frame}
Let $U$ be unitary on a complex Hilbert space, $\|e\|=1$, and put
$B_m=1+2\sum_{j=1}^{m-1}|\langle e,U^je\rangle|$.
For every integer $N$, coefficients $a_0,\ldots,a_{m-1}$ and vector
$b$,
\begin{align}
 \left\|\sum_{j=0}^{m-1}a_jU^{N+j}e\right\|^2
                  &\le B_m\sum_{j=0}^{m-1}|a_j|^2,
                         \label{eq:orbit-synthesis}\\
 \sum_{j=0}^{m-1}|\langle b,U^{N+j}e\rangle|^2
                  &\le B_m\|b\|^2.
                         \label{eq:orbit-analysis}
\end{align}
Both assertions remain true after multiplying each orbit vector by an
arbitrary scalar of modulus one.
\end{lemma}
\begin{proof}
The Gram entry at $(i,j)$ is $\langle e,U^{j-i}e\rangle$; it is
independent of $N$. Every absolute row sum is at most $B_m$.
For completeness, expansion of the squared norm and
$2|a_i||a_j|\le|a_i|^2+|a_j|^2$ bound that quadratic form by
$B_m\sum_j|a_j|^2$. This proves \eqref{eq:orbit-synthesis}.
The synthesis operator from $\mathbb C^m$ to the Hilbert space has
norm at most $\sqrt{B_m}$; its adjoint has the same norm and is the
analysis map in \eqref{eq:orbit-analysis}. Multiplication by unit
scalars conjugates the Gram matrix by a diagonal unitary. This proves
the final assertion without assuming that the vectors are orthogonal
or linearly independent.
\end{proof}

\begin{theorem}[Uncompressed moving averages]
\label{thm:direct-moving-averages}
Let $m=m(r)$, $0<|z|=r\le r_0$, and $T\ge m$ be an integer.
For all integers $N$, every $b\in\mathcal H_R$ and every
$\xi\in\R$,
\begin{align}
 \frac1T\sum_{j=0}^{T-1}
     |\langle b,\mathcal U_z^{N+j}e\rangle_R|^2
                         &\le Cr^\kappa\|b\|_2^2,
                            \label{eq:direct-moving-analysis}\\
 \left\|\frac1T\sum_{j=0}^{T-1}
      e^{-ij\xi}\mathcal U_z^{N+j}e\right\|_2^2
                         &\le Cr^\kappa.
                            \label{eq:direct-moving-synthesis}
\end{align}
The constants are independent of $N,T,R,\xi$ and $b$. In particular
these estimates apply to the original characteristic functions at
arbitrarily late return counts, not to powers of a compressed operator.
\end{theorem}
\begin{proof}
Divide the $T$ indices into complete blocks of length $m$ and one
possibly shorter block. The estimate for a shorter block follows from
the same Gram argument with its row sum bounded by $\mathcal B_m$.
There are at most $T/m+1\le2T/m$ blocks. Summing
\eqref{eq:orbit-analysis} proves
\eqref{eq:direct-moving-analysis} with constant
$2\mathcal B_m/m$.

For synthesis, a complete block of $m$ unit-modulus coefficients has
norm at most $\sqrt{m\mathcal B_m}$, and a shorter block has no
larger norm. The triangle inequality across the at most $2T/m$
blocks, followed by division by $T$, gives squared norm at most
$4\mathcal B_m/m$. Proposition~\ref{prop:short-lag-row-budget}
finishes both bounds. Modulation by $e^{-ij\xi}$ does not change any
absolute Gram entry. Thus there is no restriction on the peripheral
angle $\xi$ in this theorem.
\end{proof}

\begin{corollary}[One feasible scale on the full small annulus]
\label{cor:direct-annular-averages}
For sufficiently large $n$ put
\begin{equation}\label{eq:direct-annulus}
 \mathcal A_n=\{z\in\R^4:2n^{-99/200}\le|z|\le n^{-2/5}\}.
\end{equation}
Its rescaled radii are $2n^{1/200}$ and $n^{1/10}$.
For every $z\in\mathcal A_n$ one has
\begin{equation}\label{eq:direct-annular-scale}
 m(|z|)\le2^{-200/99}n<n/4,
 \qquad |z|^\kappa\le n^{-4/495}.
\end{equation}
Consequently both conclusions of
Theorem~\ref{thm:direct-moving-averages} hold for $T\ge n$ with
$Cr^\kappa$ replaced by $Cn^{-4/495}$, uniformly throughout the
annulus and at every peripheral angle.
\end{corollary}
\begin{proof}
The first inequality follows by raising the lower physical radius to
power $-200/99$; the exponent product is exactly one. The second follows
from the upper radius and $(2/5)(2/99)=4/495$. Thus a single explicit
choice of the lag window is compatible with every averaging time
$T\ge n$. No spatial mesh or long-orbit BV budget enters this choice.
\end{proof}

The use of the word ``average'' in these statements is essential.
They bound a sum of squared matrix coefficients or the norm of a
Cesaro average of the actual orbit. They do not bound
$|c_{n,R}(z)|$ at every specified $n$, and they do not assert that
$\|\mathcal U_z^n\|$ decays; that norm equals one. The next section
records the weighted and peripheral consequences at precisely this
level.
